\documentclass[12pt]{article}
\usepackage[margin=0.85in]{geometry}                
\geometry{letterpaper}                  
\usepackage{graphicx}
\usepackage{booktabs} % For professional looking tables
\usepackage{multirow}
\usepackage{array}
\usepackage{amsmath, amssymb, amsthm}
\usepackage[most]{tcolorbox}
\usepackage{hyperref, multicol, booktabs, xcolor, dsfont, tcolorbox}
\usepackage{tabularx}

\hypersetup{colorlinks=false, allcolors=blue}
\newcommand{\noin}{\noindent}          

\newcommand{\Bern}{\textnormal{Bern}}
\newcommand{\Bin}{\textnormal{Bin}}                      
\newcommand{\HGeom}{\textnormal{HGeom}} 
\newcommand{\Geom}{\textnormal{Geom}}   
\newcommand{\NBin}{\textnormal{NBin}}       
\newcommand{\Pois}{\textnormal{Pois}}        
\newcommand{\Unif}{\textnormal{Unif}} 
\newcommand{\Expo}{\textnormal{Expo}}
\newcommand{\N}{\mathcal{N}}     

\newcommand{\Var}{\textnormal{Var}}      
\newcommand{\Exp}{\textnormal{E}} 
\newcommand{\Cov}{\textnormal{Cov}}
\newcommand{\E}{\mathbb{E}}
\newcommand{\Prob}{\mathbb{P}}
\usepackage{graphicx} % Required for inserting images


\begin{document}
\begin{titlepage}
    \vspace*{\fill}
    \begin{center}
        {\Huge \textbf{Notes for Stat 111: Introduction to Statistical Inference} \\[0.6cm]}
        {\Large Eddy Kang \\[0.4cm]}
        {\large Spring 2026}
    \end{center}
    \vspace*{\fill}
\end{titlepage}
\newpage
\newpage

\section{Introduction}

\section{Models and Estimation}

\section{Asymptotics}

\section{Confidence Intervals}

\section{Regression}

\section{Exponential Families and Data Compression}

\section{Hypothesis testing}

\subsection{Fundamentals of hypothesis testing}
We'll now cover hypothesis testing. Say we've collected data $\textbf{Y} = \textbf{y}$, and our data is generated by $F_{Y;\theta}$. Then we define a statistical hypothesis below:

\medskip

\begin{tcolorbox}[breakable, title=\textbf{Definition}: Statistical hypothesis]
    
    We partition the parameter $\Theta$ into two disjoint sets $\Theta_0$ and $\Theta_1$, such that $\Theta_0 \cup \Theta_1 = \emptyset$. Then, we define the \underline{null} ($H_0$) and \underline{alternative} ($H_1$) hypotheses below:
    \begin{align*}
        H_0 : \theta \in \Theta_0, \ H_1 : \theta \in \Theta_1
    \end{align*}
    We call a hypothesis \underline{simple} if it is simply a point (ex: $\Theta_0 = \{\theta_0\})$, and \underline{composite} otherwise.
\end{tcolorbox}

\medskip

\noindent We want a procedure to carry out to test whether we can reject or retain a hypothesis based on the data $\textbf{y}$. Thus, we define a procedure known as a hypothesis test below.

\medskip

\begin{tcolorbox} [breakable, title=\textbf{Definition}: Statistical hypothesis test]
    
    We define a procedure to either reject or retain $H_0$ based on whether the data $\textbf{y}$ falls within a specified region.
    \begin{align*}
        \text{reject } H_0 : \textbf{y} \in A^c, \text{ the \underline{rejection region}}\\
        \text{retain } H_0 : \textbf{y} \in A, \text{ the \underline{retention region}}
    \end{align*}

    Often, it's easier to base our test on a test-statistic $T(\textbf{y})$ computed from our data. In this, case, we define a \underline{critical region} $D$ such that the following is true.
    \begin{align*}
        \text{reject } H_0: T(\textbf{y}) \in D
    \end{align*}
    The set $D$ can take multiple forms such as $D = \{T(\textbf{y}) >c\}$ or $\{T(\textbf{y} ) : c_L < T(\textbf{y}) \text{ or } T(\textbf{y}) <c_U\}$, where $c, c_L,c_U$ are the \underline{critical values}.

\end{tcolorbox}

\medskip

\noindent As we see above, we like to work with a test-statistic computed from the data, then see whether the test-statistic falls within a certain region. After we've performed the hypothesis test, we either reject or retain the null. Note that a hypothesis test compresses our data into a binary outcome. We don't know the truth, so it's possible that we make errors in our binary decision, which are described below.

\medskip

\begin{tcolorbox}[breakable, title=\textbf{Definition}: Errors in hypothesis testing]
    (1) \underline{Type I error}: We reject $H_0$, but $H_0$ is true ("\underline{false positive}").\\
    (2) \underline{Type II error}: We retain $H_0$, but $H_1$ is true ("\underline{false negative}").
\end{tcolorbox}
    
\medskip

\noindent We want to be able to evaluate our hypothesis tests in terms of the probability that we make the kinds of errors specified above. This will tell how ''good`` our test is, and how likely it is that we made the wrong judgement call. We define the power function below.

\medskip

\begin{tcolorbox}[breakable, title=\textbf{Definition}: Power]
    Assume the data were generated $F_{\textbf{Y};\theta}$, and we have our hypothesis and retention region $A$. The power function of the test is
    \begin{align*}
        \beta(\theta) = \Prob_\textbf{Y}(\textbf{Y}\notin A; \theta) = \int \mathds{1}(\textbf{y}\in A^c)f_\textbf{Y}(\textbf{y};\theta)d\textbf{y}
    \end{align*}

    Intuitively, we can think of $\beta(\theta)$ as the ``probability of rejecting the null under a specified value of $\theta$.''
\end{tcolorbox}

\medskip

\noindent Lastly, we'll define the size of a test, which will also help us evaluate a given testing procedure.

\medskip

\begin{tcolorbox}[breakable, title = \textbf{Definition}: Size]
    The size $\alpha$ of a testing procedure is defined as follows:
    \begin{align*}
        \alpha = \underset{\theta \in \Theta_0}{\max} \beta(\theta)
    \end{align*}
    Intuitively, we can think of $\alpha$, as the the ``maximal probability of rejecting the null, when in fact the null is true.'' It can also be thought of as a ``worst-case false-positive rate.'' Typically, the researcher would set $\alpha$ in advance before performing the test.
\end{tcolorbox}
\medskip

\noindent For an ideal test, we want the following to be true:
\begin{align*}
    \beta(\theta) \text{ low for } \theta \in \Theta_0, \ \beta(\theta) \text{ high for } \theta \in \Theta_1, \ \alpha \text{ low}
\end{align*}
This would mean that we reject the null often when the alternative is true, and we reject the null infrequently when the null is true. Finally, we have a small maximal probability of rejecting the null when in fact it's true.

\medskip

\noindent Now, let's outline the typical framework for frequentist hypothesis testing.
\begin{enumerate}
  \item State $H_0, H_1$ for data $\textbf{y}$.
  \item State size $\alpha$ and test-statistic $T(\textbf{Y})$.
  \item Determine critical values to determine critical region.
  \item Make decision to reject or retain $H_0$ based on value of $T(\textbf{y})$ computed from the data.
\end{enumerate}

\subsection{Carrying out hypothesis tests}

\subsection{Controlling test size}

\subsection{p-values}

\subsection{Likelihood-based tests}

\section{Bayesian statistics}

\subsection{Bayesian thinking}
In the frequentist framework, we treated $\theta$ as a fixed value, not a random variable. This meant that we had to a lot of bending-over-backwards to avoid making probability statements about $\theta$. For example, when we defined a CI for $\theta$, the interval itself was random! Under the Bayesian framework we will treat $\theta$ as an r.v., then write down a "full probability model." 

\medskip

\noindent We'll work with a parametric model $F_{\textbf{Y}|\theta}$ for data $\textbf{Y}$. Note the conditioning bar, which we can use instead of ; since $\theta$ is now an r.v. Bayesian thinking allows us to default back to probability theory, following the two principles below:
\begin{enumerate}
    \item Always obey probabiity laws
    \item Model all uncertainty using probability
\end{enumerate}

\noindent Bayesian statistics is based on Bayes' rule. The idea is to use Bayes' rule to obtain the posterior distribution of $\theta$, the belief we have in $\theta$ after having seen the data $\textbf{y}$.

\medskip

\begin{tcolorbox}[breakable, title = \textbf{Definition}: Posterior distribution]
    We apply Bayes' rule to obtain the \underline{posterior} distribution of $\theta$, denoted $\pi(\theta|\textbf{y})$.
    \begin{align*}
        \pi(\theta |\textbf{y}) = \dfrac{f(\textbf{y}|\theta)\pi(\theta)}{f(\textbf{y})} \propto \mathcal{L}(\theta;\textbf{y})\pi(\theta)
    \end{align*}
    Note that since the \underline{marginal likelihood} $f(\textbf{y})$ is a normalizing constant, we can write the posterior as proportional to the \underline{likelihood} times the \underline{prior}, which quantitifes our belief in $\theta$ prior to seeing the data. The marginal likelihood can always be calculated:
    \begin{align*}
        f(\textbf{y}) = \int_{-\infty}^{\infty} f(\textbf{y}|\theta)\pi(\theta)d\theta
    \end{align*}
    
\end{tcolorbox}
\medskip

\noindent What's nice about the Bayesian approach is that there's no controversy to our calculation above. Once we've defined a prior $\pi(\theta)$ and a model $f_{\textbf{Y}|\theta} (\textbf{y})$, we simply apply the laws of probability. The caveat is that writing down a prior and a model is a high barrier of entry, and is difficult in practice.

\medskip
\noindent Let's take a look at an example of how Bayes' rule can be applied to obtain the posterior distribution $\theta |\textbf{y}$.

\medskip

\begin{tcolorbox}[breakable, title = \textbf{Example 9.2.3 in textbook}: Single Bernoulli trial with Beta prior]
    Assume that our data is $Y_1$. Our model is $Y_1|\theta \sim \Bern(\theta)$, and our prior is $\theta \sim \text{Beta}(\alpha,\beta)$. Let's write down the likelihood function:
    \begin{align*}
        \mathcal{L}(\theta;Y_1) = \Prob(Y_1 = y_1|\theta)= \theta^{y_1}(1-\theta)^{1-y_1}
    \end{align*}
    Thus, applying the Beta PDF and Bayes' rule, we can write down the posterior distribution $\pi(\theta|\textbf{y})$.
    \begin{align*}
        \pi(\theta|\textbf{y}) 
        &\propto \mathcal{L}(\theta;Y_1)\pi(\theta) \\
        &\propto \left\{\theta^{y_1}(1-\theta)^{1-y_1}\right\}\left\{\theta^{\alpha-1}(1-\theta)^{\beta-1}\right\} \text{ by Beta PDF}\\
        &= \theta^{y_1 +\alpha -1}(1-\theta)^{1-y_1+\beta+1}
    \end{align*}
    This is simply the $\text{Beta}(\alpha +y_1, \beta + 1 - y_1)$ PDF, without the normalizing constant! Thus, we have the following result.
    \begin{align*}
        \boxed{\theta|\textbf{y}\sim \text{Beta}(\alpha + y_1, \beta + 1 - y_1)}
    \end{align*}
    Note that this result is an example of \underline{Beta-Binomial} conjugacy! When a Beta prior is fed through a Binomial model, we still get a Beta posterior.
\end{tcolorbox}

\medskip

\noindent We can actually generalize the result above to $n$ Bernoulli datapoints:

\medskip

\begin{tcolorbox}[breakable, title= \textbf{Example 9.2.3 in textbook (cont.)}: $n$ ind. Bernoulli trials with Beta prior]
    Assume our data is $\textbf{Y} = Y_1,\dots,Y_n$, where $Y_j |\theta\overset{\text{i.i.d.}}{\sim} \Bern(\theta)$, with prior $\theta \sim \text{Beta}(\alpha, \beta)$.
    We have the likelihood function below:
    \begin{align*}
        \mathcal{L}(\theta, \textbf{Y}) 
        &= \prod_{j=1}^{n}f_{Y_1|\theta}(y_1) \ \text{by ind.}\\
        &= \theta^{n\bar{y}}(1-\theta)^{n(1-\bar{y})} \ \text{by Bern. PMF}
    \end{align*}
    Plugging in this likelihood into the expression for $\pi(\theta|\textbf{y})$, we obtain the result
    \begin{align*}
        \pi(\theta|\textbf{y}) = \theta^{\alpha+n\bar{y}-1 }(1-\theta)^{\beta + n(1-\bar{y}) - 1}
    \end{align*}
    Thus, by the Beta PDF,
    \begin{align*}
        \boxed{\theta|\textbf{y} \sim \text{Beta}(\alpha +n\bar{y}, \beta + n(1-\bar{y}))}
    \end{align*}
\end{tcolorbox}

\medskip

\noindent Priors must be constructed carefully. Note the following:

\medskip

\begin{tcolorbox}[breakable, title = \textbf{Remark}: Cromwell's rule]
    ``When constructing a prior, don't set $\pi(\theta)$ for any value of $\theta$!'' This is because if we let the prior be 0 anywhere, the posterior $\theta|\textbf{y}$ will always be 0 no matter how much the data might suggest that that value of $\theta$ is a good fit for the data.
\end{tcolorbox}

\subsection{Bayesian point estimators}

\noindent The ``output'' of the Bayesian workflow is the posterior distribution $\pi(\theta|\textbf{y})$. But, it's often easier to work with point values rather than the entire posterior distribution. In other words, we want a clear summary of our output. Some common Bayesian point estimators are listed below:

\medskip

\begin{tcolorbox}[breakable, title = \textbf{Definition}: \text{Posterior mean, median, mode}]
    Suppose we've obtained the posterior distribution $\pi(\theta|\textbf{y})$. We define the following quantities.
    \begin{align*}
        &1) \text{ \underline{Posterior mean}: }\E[\theta|\textbf{y}]  = \int_{\infty}^{\infty}\theta\pi(\theta|\textbf{y})d\theta\\
        &2) \text{ \underline{Posterior median}: }
        Q_{\theta|\textbf{y}}(0.5)\\
        &3) \text{ \underline{Posterior mode}: }
        \underset{\theta \in \Theta}{\text{argmax }}\pi(\theta|\textbf{y})
    \end{align*}
    The posterior mode is also called the \underline{maximum \emph{a posteriori} (MAP) estimator}.
\end{tcolorbox}

\medskip
\noindent Often, we'll define a Bayesian estimator $\hat\theta$ as a quantity that minimizes the posterior expected loss function, a function we defined to measure the distance between our true $\theta$ and our estimate $\hat\theta$.

\medskip
\begin{tcolorbox}[breakable, title = \textbf{Definition}: Posterior expected loss]
    We construct a general Bayesian estimator $\hat\theta_\text{Bayes}$ using a \underline{loss function} Loss$(\theta,\hat\theta)$. 
    \begin{align*}
        \hat\theta_\text{Bayes}= \underset{\hat\theta \in \Theta}{\text{argmin }}\E[\text{Loss}(\theta,\hat\theta)|\textbf{y}]
    \end{align*}
    Our estimator minimizes the posterior expected loss: the expected loss given the data $\textbf{y}$, which can be computed using the posterior distribution $\theta|\textbf{y}$.
\end{tcolorbox}

\medskip

\noindent We can actually derive the posterior mean and mode using the definition above.

\medskip
\begin{tcolorbox}[breakable, title = \textbf{Theorem 9.3.3 in textbook}: \text{Origins of posterior mean, median}]
    If we use the \underline{squared-error loss}, $\text{Loss}(\theta,\hat\theta) =(\theta-\hat\theta)^2$, then the posterior mean minimizes the posterior expected loss:
    \begin{align*}
        \E[\theta|\textbf{y}] = \underset{c \in \Theta}{\text{argmin }}\E[(\theta-c)^2|\textbf{y}]
    \end{align*}
    If we used the \underline{absolute error loss}, $\text{Loss}(\theta,\hat\theta) = |\theta-\hat\theta|$, then the posterior median minimizes the posterior expected loss:
    \begin{align*}
        Q_{\theta|\textbf{y}}(0.5) = \underset{c \in \Theta}{\text{argmin }}\E[|\theta-c||\textbf{y}]
    \end{align*}
\end{tcolorbox}

\medskip
\noindent Using some clever tricks, we can also derive a statement about the posterior mean for Gaussian data that generalizes to all possible priors.

\medskip

\begin{tcolorbox}[breakable, title = \textbf{Example 9.3.4 in textbook}: Tweedie's formula]
    Assume Gaussian data $Y$, such that
    \begin{align*}
        f_{Y|\theta}(y)  = \dfrac{1}{\sqrt{2\pi\sigma^2}}\exp\left\{-\dfrac{1}{2\sigma^2}(y-\theta)^2\right\}
    \end{align*}
    Then for any prior on $\theta$, the following statement holds:
    \begin{align*}
        \boxed{\E[\theta|y]= y + \sigma^2 \dfrac{\partial\log f(y)}{\partial y}}
    \end{align*}
\end{tcolorbox}

\medskip

\noindent Tweedie's formula is fascinating, because it means that for Gaussian data, we don't need any information about $\theta$ to obtain the posterior mean! 

\medskip

\noindent Next, let's explore some properties of the MAP estimator. This estimator is convenient because we don't ever need to compute the normalizing constant. This is because we don't need the posterior distribution $\theta|\textbf{y}$. We can simply compute the MAP estimate using the likelihood and prior.

\medskip

\begin{tcolorbox}[breakable, title = \textbf{Remark}: Properties of the MAP estimator] 
    Recall the definition of the MAP estimator:
    \begin{align*}
        \hat\theta_\text{MAP} = \underset{\theta \in \Theta}{\text{argmax }}\pi(\theta|\textbf{y})
    \end{align*}
    We can express the posterior distribution in terms of the likelihood, prior, and some constant $c$.
    \begin{align*}
        \pi(\theta|\textbf{y}) \propto \mathcal{L}(\theta;\textbf{y})\pi(\theta) \Rightarrow\log\pi(\theta|\textbf{y}) = c + \log \mathcal{L}(\theta;\textbf{y}) + \log \pi(\theta)
    \end{align*}
    Thus, we can rewrite the MAP in terms of only the likelihood and prior:
    \begin{align*}
        \hat\theta_\text{MAP} 
        &= \underset{\theta \in \Theta}{\text{argmax }}\log \pi(\theta|\textbf{y})\\
        &=  \underset{\theta \in \Theta}{\text{argmax }} \left\{c + \log \mathcal{L}(\theta;\textbf{y}) + \log \pi(\theta)\right\}\\
        &= \underset{\theta \in \Theta}{\text{argmax }} \left\{\log \mathcal{L}(\theta;\textbf{y}) + \log \pi(\theta)\right\}
    \end{align*}
    Note that if our prior is Uniform (constant with respect to $\theta) $, then the MAP is simply the MLE!
\end{tcolorbox}

\subsection{Credible intervals}

\noindent In frequentist statistics, we had the confidence interval. For Bayesian statistics, we have the credible interval. 

\medskip

\begin{tcolorbox}[breakable, title= \textbf{Definition}: Credible interval]
    We define a $1-\alpha$ credible interval for $\theta$ using the constants $L(\textbf{y})$ and $U(\textbf{y})$.
    \begin{align*}
        \Prob\left\{L(\textbf{y}) \leq \theta \leq U(\textbf{y})|\textbf{y}\right\} = 1 -\alpha
    \end{align*}
\end{tcolorbox}

\medskip

\noindent The credible interval is exactly the opposite of a confidence interval. In a confidence interval, $\theta$ was fixed and our bounds $L(\textbf{Y})$ and $U(\textbf{Y})$ were r.v.s. In our credible interval,  $\theta$ is the r.v. and $L(\textbf{y})$ and $U(\textbf{y})$ are constants. In the Bayesian lens, we can make a direct probability statement about where $\theta$ is located!

\medskip

\noindent We can define a credible interval using the quantile function as we previously did for a confidence interval.

\begin{tcolorbox}[breakable, title = \textbf{Remark}: Defining credible interval with quantiles]
    In general, for a $1-\alpha$ credible interval, we can define the interval
    \begin{align*}
        [Q_{\theta|\textbf{y}}(\alpha/2), Q_{\theta|\textbf{y}}(1- \alpha/2)]
    \end{align*}
    This is equivalent to us taking a central "slice" of the posterior distribution $\pi(\theta|\textbf{y})$, for which the area under the curve is $1-\alpha$.
\end{tcolorbox}

\medskip

\subsection{Computing Bayesian estimators via simulation}

\noindent Although Bayesian estimators have nice probablisitic properties and are straightforward to analyze in theory, they are difficult to actually simulate. They often involve hard integrals like the following:

\medskip

\begin{tcolorbox}[breakable, title= \textbf{Example}: Computing the posterior mean]
Calculating Bayesian estimators often leads to computational challenges, such as very hard integrals!
    \begin{align*}
        \E[\theta|\textbf{y}] = \int\theta\pi(\theta|\textbf{y})d\theta = \dfrac{\int\theta \pi(\theta)\mathcal{L}(\theta;\textbf{y})d\theta}{\int\pi(\theta)\mathcal{L}(\theta;\textbf{y})d\theta}
    \end{align*}
\end{tcolorbox}

\medskip

\noindent We need a better way to do this. Simulation seems promising. We know that Markov-chain Monte Carlo (MCMC) can simulate draws from a given continuous distribution $f$. We can apply this to our posterior distribution to obtain estimates of our Bayesian estimators.

\medskip

\begin{tcolorbox}[breakable, title= \textbf{Example}: Bayesian simulation]
    Say have the posterior distribution $\pi(\theta|\textbf{y})$. We want to compute 2 Bayesian estimators: the posterior mean $\E[\theta|\textbf{y}]$ and the posterior median $Q_{\theta|\textbf{y}}(0.5)$. However, these integrals are too hard! We simply sample draws from $\pi(\theta|\textbf{y})$ using MCMC:
    \begin{align*}
        \theta_1^*, \dots,\theta_B^* \sim \theta|y
    \end{align*}
    Then we estimate the posterior mean and posterior median from the data!
    \begin{align*}
        \E[\theta|\textbf{y}] \approx \dfrac{1}{B}\sum_{b=1}^{B}\theta_b^*, \ Q_{\theta|\textbf{y}}(0.5) \approx \hat Q_{\lceil 0.5B\rceil}(0.5)
    \end{align*}
\end{tcolorbox}

\medskip

\noindent Note that this strategy is invariant under reparameterization of $\theta$.

\medskip

\begin{tcolorbox}[breakable, title= \textbf{Example}: Bayesian simulation using reparameterization]
    Say our friend wants to calculated the mean posterior mean and median of $\lambda$ instead of $\theta$, where $\lambda = h(\theta)$. Then, we simply apply the transformation $h$ to our simulated draws from $\pi(\theta|\textbf{y})$.
    \begin{align*}
        \lambda_b^* = h(\theta_b^*) \text{ for }b=1,\dots, B
    \end{align*}
    Then, we again just estimate the posterior mean and posterior median from the data.
    \begin{align*}
        \E[\lambda|\textbf{y}] \approx \dfrac{1}{B}\sum_{b=1}^{B}h(\theta_b^*), \ Q_{\lambda|\textbf{y}}(0.5) \approx \hat Q_{\lceil 0.5B\rceil}(0.5)
    \end{align*}
    We can do all of this without worrying about change of variables or the Jacobian, which is very convenient.
    
\end{tcolorbox}

\medskip

\noindent Simulation is especially powerful when our parameter $\theta$ is multivariate: $\theta = (\psi, \lambda)$. Suppose we want to study $\lambda$ alone.

\begin{tcolorbox}[breakable, title= \textbf{Example}: Bayesian simulation using marginalization]
    Suppose $\theta = (\psi,\lambda)$, and we've simulated draws from $\pi(\theta|\textbf{y})$:
    \begin{align*}
        \theta_1^*,\dots,\theta_B^* = \binom{\psi_1^*}{\lambda_1^*} ,\dots,\binom{\psi_B^*}{\lambda_B^*} \sim \theta|\textbf{y}
    \end{align*}
    We would like to study $\lambda$, so we want $\pi(\lambda|\textbf{y}) = \int\pi(\psi,\lambda|\textbf{y})d\textbf{y}$ by marginalization. This is very hard to compute. So, instead, we marginalize by just throwing out the $\psi_b^*$. We're left with
    \begin{align*}
        \lambda_1^*,\dots,\lambda_B^* \sim \lambda|\textbf{y}
    \end{align*}
    Now, we can estimate quantities such as $\E[\lambda|\textbf{y}]$ or $Q_{\lambda|\textbf{y}}(0.5)$ just as we did above.
\end{tcolorbox}

\subsection{Conjugate priors}

\noindent It turns out that some priors and likelihoods work especially well together if they have a particular distribution. We've seen conjugacy in Stat 110 already, so let's explore it deeper.

\medskip

\begin{tcolorbox}[breakable, title = \textbf{Definition}: Conjugate prior]
    Recall Bayes' rule:
    \begin{align*}
        \pi(\theta|\textbf{y}) \propto \pi(\theta)\mathcal{L}(\theta;\textbf{y})
    \end{align*}
    A family of priors $\pi(\theta)$ is \underline{conjugate} for a statistical model if choosing a prior in that family leads to a posterior also in the same family.
\end{tcolorbox}

\medskip

\noindent We've seen a couple examples of this in the past.

\medskip

\begin{tcolorbox}[breakable, title = \textbf{Example}: Beta-Binomial and Gamma-Poisson conjugacy]
    Beta-Binomial conjugacy:
    \begin{align*}
        Y|p \sim \Bin(n,p), \ p \sim \text{Beta}(a,b) \Rightarrow p|(Y=y) \sim \text{Beta}(a+y, b+n-y)
    \end{align*}
    Gamma-Poisson conjugacy:
    \begin{align*}
        Y|\lambda\sim \Pois(\lambda), \ \lambda \sim \text{Gamma}(r_0, b_0) \Rightarrow \lambda|(Y=y) \sim \text{Gamma}(r_0+y, b_0+1)
    \end{align*}
\end{tcolorbox}

\medskip

\noindent Next, we extend this to examine conjugacy involving the Normal.

\medskip

\begin{tcolorbox}[breakable, title = \textbf{Example 9.5.4 in textbook}: Normal-Normal conjugacy (single observation)]
    Suppose we have a single observation $Y$. We define $Y|\mu\sim \N(\mu,\sigma^2)$ and $\mu \sim \N(m_0, \tau_0^2)$. Assume $\sigma^2, m_0, \tau_0^2$ known. We want to find the posterior $\mu|y$. By Bayes' rule,
    \begin{align*}
        \pi(\mu|y) = c\cdot \pi(\mu)\mathcal{L}(\mu;y) \ \text{where $c$ is a constant}
    \end{align*}
    We examine the log posterior,
    \begin{align*}
        \log \pi(\mu|y) 
        &= c + \log \pi(\mu) + \log \mathcal{L}(\theta;y)\\
        &= c'-\dfrac{1}{2\tau_0^2}(\mu-m_0)^2 + \dfrac{1}{2\sigma^2}(y-\mu)^2\\
        &= c'' -\dfrac{\mu^2}{2}\left(\dfrac{1}{\sigma^2}-\dfrac{1}{\tau_0^2}\right) + \mu\left(\dfrac{y}{\sigma^2} + \dfrac{m_0}{\tau_0^2}\right)\\
        &= c'' - \dfrac{\mu^2}{2\tau_1^2}+\dfrac{\mu}{\sigma_1^2}m_1
        \end{align*}
        Here, we've made the following substitutions:
        \begin{align*}
            \tau_1^{-2} = \sigma^{-2} + \tau_0^{-2}, \ m_1 = \tau_1^2\left(\dfrac{y}{\sigma^2} + \dfrac{m_0}{\tau_0^2}\right)
        \end{align*}
        Thus, pattern-matching to the log density of the Normal,
        \begin{align*}
            \boxed{\mu|y \sim \N(m_1,\tau_1^2)}
        \end{align*}
        Note that the $\sigma^{-2}$ and $\tau_0^2$ are the \underline{precisions}, meaning that the precisions of the likelihood and prior add!
\end{tcolorbox}

\medskip
We can utilize the same logic to examine Normal-Normal conjugacy for $n$ datapoints.

\medskip

\medskip
\begin{tcolorbox}[breakable, title = \textbf{Theorem 9.5.5 in textbook}: Normal-normal conjugacy ($n$ observations)]
    We generalize the above to $n$ observations. Suppose we have data $Y_1,\dots,Y_n$, where $Y_j|\mu \overset{\text{i.i.d.}}{\sim} \N(\mu,\sigma^2)$ and $\mu \sim \N(m_0,\tau_0^2)$. Then, 
\end{tcolorbox}
\subsection{Predictive posterior}

\subsection{Hierarchical models}
\noindent 

\section{Sampling and resampling}

\section{Causal inference}
\end{document}
